\chapter{The .X16S File Format}
\label{app:fileformat}

This appendix describes \ccprod's native workbook format precisely enough to
write a reader for it. Nothing here is needed to use the program.

All integers are little-endian, fixed width, with no alignment and no padding.

\section{The envelope}

\begin{ccfigure}
\resizebox{\linewidth}{!}{%
\begin{tikzpicture}[x=1pt,y=-1pt,
  fld/.style={draw=ccink,line width=0.6pt,minimum height=17pt,inner sep=2pt,
              font=\ttfamily\fontsize{7.6}{9}\selectfont,anchor=west},
  hd/.style={fld,fill=black!12},
  ann/.style={font=\sffamily\fontsize{7}{8.6}\selectfont,text=ccink,
              anchor=north}]
  \def\y{0}
  \node[hd,minimum width=44pt]  (a) at (0,\y)   {"X16S"};
  \node[hd,minimum width=40pt]  (b) at (44,\y)  {ver};
  \node[hd,minimum width=40pt]  (c) at (84,\y)  {flags};
  \node[fld,minimum width=48pt] (d) at (124,\y) {SHNM};
  \node[fld,minimum width=44pt] (e) at (172,\y) {STYL};
  \node[fld,minimum width=44pt] (f) at (216,\y) {STRS};
  \node[fld,minimum width=44pt] (g) at (260,\y) {COLW};
  \node[fld,minimum width=42pt] (h) at (304,\y) {CELS};
  \node[fld,minimum width=38pt] (i) at (346,\y) {META};
  \node[hd,minimum width=38pt]  (j) at (384,\y) {"END "};
  \node[hd,minimum width=44pt]  (k) at (422,\y) {CRC-32};

  \node[ann] at (22,12)  {4};
  \node[ann] at (64,12)  {2};
  \node[ann] at (104,12) {2};
  \node[ann] at (148,12) {chunk};
  \node[ann] at (194,12) {chunk};
  \node[ann] at (238,12) {chunk};
  \node[ann] at (282,12) {chunk};
  \node[ann] at (325,12) {chunk};
  \node[ann] at (365,12) {chunk};
  \node[ann] at (403,12) {chunk};
  \node[ann] at (444,12) {4};

  \draw[ccink,line width=0.5pt] (0,34) -- (0,40) -- (422,40) -- (422,34);
  \node[font=\sffamily\fontsize{7}{8.6}\selectfont,text=ccink,anchor=north]
    at (211,41) {the CRC is computed over everything before it};
\end{tikzpicture}}
\caption{The layout of a \fname{.X16S} file}
\label{fig:x16slayout}
\end{ccfigure}

\begin{cctablethree}{Offset}{Size}{Field}{0.75in}{0.6in}{3.45in}
0 & 4 & Magic, the four bytes \lit{X16S} \\
4 & 2 & Version. Currently 3. A higher number is refused with
\msg{Damaged or unrecognised file}. \\
6 & 2 & Flags. Always written as 0; read and discarded. \\
8 & \ldots & The chunks, in the order shown above \\
\ldots & 8 & The terminating chunk: tag \lit{END\ } and length 0 \\
\ldots & 4 & CRC-32 of every byte before this field \\
\end{cctablethree}

\section{Chunks}

Every chunk has the same three-part shape:

\begin{cctablethree}{Offset}{Size}{Field}{0.75in}{0.6in}{3.45in}
0 & 4 & Four-character tag \\
4 & 4 & Length of the payload, not counting these eight bytes \\
8 & \emph{length} & Payload \\
\end{cctablethree}

\textbf{A reader must skip a chunk whose tag it does not know} rather than
rejecting the file. That is the point of the framing, and it is what lets a
later version of \ccprod{} add to the format without breaking this one.

Chunks are written in a fixed order, and a reader may rely on \lit{STRS}
preceding \lit{CELS}, because a text cell names its string by an identifier
assigned in \lit{STRS}.

\subsection{SHNM --- worksheet names}

The names, NUL-terminated, one after another, in worksheet order. Each is at
most 31 characters plus its terminator.

\subsection{STYL --- number formats}

\begin{cctablethree}{Offset}{Size}{Field}{0.75in}{0.6in}{3.45in}
0 & 2 & Count of style records \\
2 & 5 each & The style records \\
\end{cctablethree}

Each record is five bytes:

\begin{cctablethree}{Offset}{Size}{Field}{0.75in}{0.6in}{3.45in}
0 & 1 & Number format: 0 General, 1 Integer, 2 Decimal, 3 Currency, 4 Percent,
5 Date, 6 Time, 7 Text \\
1 & 1 & Foreground colour \\
2 & 1 & Background colour \\
3 & 1 & Flags: bit 0 bold; bits 1--2 alignment (0 automatic, 1 left, 2 right,
3 centre) \\
4 & 1 & Decimal places \\
\end{cctablethree}

The chunk length must be exactly $2 + 5n$ or the file is rejected. Style~0 is
written but not read back: the reader has already created a General style at
index 0 and skips the file's first record.

\subsection{STRS --- the text pool}

\begin{cctablethree}{Offset}{Size}{Field}{0.75in}{0.6in}{3.45in}
0 & 2 & Count of string records \\
2 & \ldots & The records \\
\end{cctablethree}

Each record is a 2-byte identifier, a 2-byte length, and that many bytes of
text. Identifiers start at 1 and are the identifiers the cells refer to. A
string longer than 1,024 bytes makes the file invalid.

\begin{ccnote}
Text is stored as plain bytes in the machine's own character set, not as
PETSCII and not as UTF-8. What is in the file is what appears on screen.
\end{ccnote}

\subsection{COLW --- column widths}

One chunk per worksheet: a 1-byte worksheet number followed by 256 width bytes.
All 256 are always written whether or not they differ from the default of ten.
A short chunk is read as far as it goes; a long one has its tail ignored.

\subsection{CELS --- the cells}

One chunk per worksheet.

\begin{cctablethree}{Offset}{Size}{Field}{0.75in}{0.6in}{3.45in}
0 & 1 & Worksheet number \\
1 & 4 & Total cell count \\
5 & \ldots & Row blocks, until the chunk ends \\
\end{cctablethree}

Each row block is a 2-byte row number, a 2-byte count of cells in that row, and
that many cell records. Rows ascend; within a row, columns ascend. A count of
0, or of more than 256, makes the file invalid.

\subsection{META --- summary}

Four 16-bit fields: worksheet count, active worksheet, highest row used,
highest column used.

\begin{ccnote}
\lit{META} is written but \textbf{not read}. The worksheet count comes from
\lit{SHNM} and the used range is recomputed as the cells are loaded. It is
recorded here because it is in the file, not because it matters.
\end{ccnote}

\section{The cell record}

Eight bytes.

\begin{ccfigure}
\begin{tikzpicture}[x=1pt,y=-1pt,
  by/.style={draw=ccink,line width=0.6pt,minimum width=34pt,minimum height=19pt,
             inner sep=1pt,font=\ttfamily\fontsize{7.6}{9}\selectfont},
  bx/.style={by,fill=black!12},
  num/.style={font=\sffamily\fontsize{6.8}{8}\selectfont,text=ccink,anchor=south}]
  \foreach \i/\t/\s in {0/col/1,1/type/1,2/style/1}{
    \node[bx] at ({\i*35+17},10) {\t};
    \node[num] at ({\i*35+17},-2) {\i};
  }
  \foreach \i in {3,4,5,6,7}{
    \node[by] at ({\i*35+17},10) {};
    \node[num] at ({\i*35+17},-2) {\i};
  }
  \draw[ccink,line width=0.5pt] (105,22) -- (105,29) -- (280,29) -- (280,22);
  \node[font=\sffamily\fontsize{7.4}{9}\selectfont,text=ccink,anchor=north]
    at (192,30) {value, five bytes, meaning set by \texttt{type}};
\end{tikzpicture}
\caption{The eight-byte cell record}
\label{fig:cellrecord}
\end{ccfigure}

\begin{cctablethree}{Offset}{Size}{Field}{0.75in}{0.6in}{3.45in}
0 & 1 & Column, 0--255 \\
1 & 1 & Type: 0 empty, 1 number, 2 text, 3 boolean, 4 date, 5 formula, 6 error \\
2 & 1 & Style, an index into \lit{STYL} \\
3 & 5 & Value \\
\end{cctablethree}

\begin{cctable}{Type}{The five value bytes hold}{1.35in}{3.47in}
Number & The number, in the 5-byte format of Appendix~\ref{app:numbers} \\
Date & The same, holding a day count \\
Boolean & The same, holding 1 or 0 \\
Text & A 2-byte string identifier, then three zero bytes \\
Formula & A 2-byte string identifier naming the formula's \emph{source text},
then three zero bytes \\
Error & The error code in the first byte \\
\end{cctable}

A blank cell has no record at all; storage is sparse.

\section{Version history}

\begin{cctable}{Version}{Meaning}{0.75in}{4.07in}
1 & One unnamed worksheet. \lit{COLW} and \lit{CELS} carry no worksheet byte.
Still opens; its formulas are dropped. \\
2 & Every worksheet. \lit{COLW} and \lit{CELS} gain the leading worksheet byte.
Still opens; its formulas are dropped. \\
3 & A formula cell stores the identifier of its source text, so formulas
survive a save and load. Current. \\
\end{cctable}

Versions 1 and 2 stored a memory address in a formula cell, which means nothing
once the file has been read into a different machine. Those cells are read and
discarded.

\section{Saving is atomic}

A save writes the whole workbook to \fname{CMDRCALC.TMP}, reads it back and
verifies the CRC, and only then deletes the previous file and renames the
temporary one over it. Any failure removes the temporary file and leaves the
previous workbook untouched. The cost is that the card must have room for two
copies at once.

\section{Loading}

The header is validated before anything in memory is touched, so a bad header
leaves the current workbook alone. Past that point the workbook is cleared
first; a failure after that point clears it again, so a damaged file leaves an
empty workbook rather than a half-built one.

Formulas are collected as the file is read and compiled after it has been
closed, which is why a large workbook pauses twice.
